\usepackage{fontspec}
\usepackage{geometry}
\usepackage{microtype}
\usepackage{xcolor}
\usepackage{amsmath,amssymb,mathtools}
\usepackage{booktabs,longtable,tabularx,array,multirow}
\usepackage{enumitem}
\usepackage{hyperref}
\usepackage{fancyhdr}
\usepackage{titlesec}
\usepackage{listings}
\usepackage[most]{tcolorbox}
\usepackage{tikz}
\usetikzlibrary{positioning,shapes.geometric,shapes.misc,arrows.meta,decorations.pathreplacing}
\usepackage{caption}

\setlength{\headheight}{16pt}
\emergencystretch=2em
\renewcommand{\contentsname}{Tartalom}

% Fontbeallitasok: ha a DejaVu fontok nincsenek telepitve,
% akkor automatikusan a MiKTeX/TeX Live alap Latin Modern fontjaira valtunk.
\IfFontExistsTF{DejaVu Serif}{\setmainfont{DejaVu Serif}}{\setmainfont{Latin Modern Roman}}
\IfFontExistsTF{DejaVu Sans}{\setsansfont{DejaVu Sans}}{\setsansfont{Latin Modern Sans}}
\IfFontExistsTF{DejaVu Sans Mono}{\setmonofont{DejaVu Sans Mono}}{\setmonofont{Latin Modern Mono}}

\geometry{left=25mm,right=25mm,top=24mm,bottom=27mm}

\hypersetup{
  colorlinks=true,
  linkcolor=blue!45!black,
  urlcolor=blue!55!black,
  pdftitle={Zarovizsga tetel},
  pdfauthor={Baba Levente - tanulasi segedanyag}
}

\definecolor{accent}{HTML}{1F4E79}
\definecolor{lightaccent}{HTML}{EAF2F8}
\definecolor{warnbg}{HTML}{FFF4E5}
\definecolor{codebg}{HTML}{F7F7F7}

\tikzset{
  flow/.style={draw, rounded corners, align=center, minimum width=21mm, minimum height=8mm, fill=lightaccent},
  pred/.style={draw, diamond, aspect=2, align=center, inner sep=1pt, fill=warnbg},
  merge/.style={draw, circle, inner sep=1.2pt, fill=black!8},
  terminator/.style={draw, rounded corners=4mm, align=center, minimum width=20mm, minimum height=8mm, fill=black!5},
  arrow/.style={->, >=stealth, thick},
  zvflowstart/.style={draw, rounded rectangle, minimum width=2.3cm, minimum height=0.7cm, align=center},
  zvflowprocess/.style={draw, rectangle, minimum width=2.7cm, minimum height=0.7cm, align=center},
  zvflowdecision/.style={draw, diamond, aspect=2.2, minimum width=2.8cm, minimum height=0.8cm, align=center}
}


\titleformat{\section}{\Large\bfseries\sffamily\color{accent}}{\thesection.}{0.6em}{}
\titleformat{\subsection}{\large\bfseries\sffamily\color{accent}}{\thesubsection.}{0.6em}{}
\titleformat{\subsubsection}{\normalsize\bfseries\sffamily\color{accent}}{\thesubsubsection.}{0.6em}{}

\setlist[itemize]{topsep=3pt,itemsep=2pt,parsep=0pt,leftmargin=1.4em}
\setlist[enumerate]{topsep=3pt,itemsep=2pt,parsep=0pt,leftmargin=1.6em}
\setlength{\parindent}{0pt}
\setlength{\parskip}{6pt}
\renewcommand{\arraystretch}{1.25}
\newcolumntype{L}[1]{>{\raggedright\arraybackslash}p{#1}}
\renewcommand{\tabularxcolumn}[1]{>{\raggedright\arraybackslash}p{#1}}

\providecommand{\ZVHEADERLEFT}{\TETELNUMBER. tétel}
\providecommand{\ZVHEADERRIGHT}{\TETELSHORTTITLE}

\pagestyle{fancy}
\fancyhf{}
\lhead{\ZVHEADERLEFT}
\rhead{\ZVHEADERRIGHT}
\cfoot{\thepage}

\lstdefinelanguage{pseudocode}{
  morekeywords={IF,THEN,ELSE,WHILE,DO,FOR,TO,DOWNTO,REPEAT,UNTIL,RETURN,CALL,INPUT,OUTPUT,AND,OR,NOT,TRUE,FALSE},
  sensitive=true,
  morecomment=[l]{//}
}

\lstdefinestyle{code}{
  basicstyle=\ttfamily\small,
  backgroundcolor=\color{codebg},
  frame=single,
  rulecolor=\color{black!15},
  columns=fullflexible,
  keepspaces=true,
  breaklines=true,
  tabsize=2,
  showstringspaces=false,
  numbers=left,
  numberstyle=\tiny\color{black!55},
  numbersep=7pt,
  xleftmargin=2.2em,
  framexleftmargin=1.7em,
  captionpos=b
}

\lstdefinestyle{pseudo}{
  style=code,
  language=pseudocode,
  keywordstyle=\bfseries\color{accent!85!black},
  commentstyle=\itshape\color{black!55}
}

\lstdefinestyle{csource}{
  style=code,
  language=C,
  keywordstyle=\bfseries\color{accent!85!black},
  commentstyle=\itshape\color{black!55},
  stringstyle=\color{orange!70!black}
}

\lstset{style=code}
\renewcommand{\lstlistingname}{Kódrészlet}
\renewcommand{\lstlistlistingname}{Kódrészletek}
\newcommand{\zvcode}[2][]{\lstinputlisting[style=code,#1]{#2}}
\newcommand{\zvpseudocode}[2][]{\lstinputlisting[style=pseudo,#1]{#2}}
\newcommand{\zvccode}[2][]{\lstinputlisting[style=csource,#1]{#2}}

\newtcolorbox{summarybox}{
  enhanced,
  colback=lightaccent,
  colframe=accent!70!black,
  arc=2mm,
  boxrule=0.5pt,
  left=2mm,right=2mm,top=1mm,bottom=1mm
}

\newtcolorbox{warningbox}{
  enhanced,
  colback=warnbg,
  colframe=orange!70!black,
  arc=2mm,
  boxrule=0.5pt,
  left=2mm,right=2mm,top=1mm,bottom=1mm
}

\newcommand{\adt}[2]{\ensuremath{#1=(#2)}}
\newcommand{\set}[1]{\ensuremath{\{#1\}}}
\newcommand{\code}[1]{\texttt{#1}}
\newcommand{\base}[2]{\ensuremath{#1_{(#2)}}}
\newcommand{\addr}{\ensuremath{@}}


\providecommand{\TETELKIIRAS}{\TETELTITLE. \TETELTOPIC}
\providecommand{\ZVAUTHOR}{Baba Levente}
\providecommand{\ZVYEAR}{2026}
\providecommand{\ZVAID}{ChatGPT segítségével}

\newcommand{\makezvtitle}{%
\begin{titlepage}
  \centering
  \vspace*{18mm}
  {\Large\sffamily Programtervező Informatikus BSc\par}
  \vspace{1mm}
  {\large\sffamily \ZVYEAR-os záróvizsga tételek\par}
  \vspace{15mm}
  {\Huge\bfseries\sffamily \TETELNUMBER. tétel\par}
  \vspace{5mm}
  {\LARGE\bfseries\sffamily \TETELTITLE\par}
  \vspace{10mm}
  \begin{summarybox}
    \textbf{Tételkiírás:} \TETELKIIRAS
  \end{summarybox}
  \vfill
  {\large Készítette: \textbf{\ZVAUTHOR}\par}
  \vspace{2mm}
  {\large \ZVAID\par}
\end{titlepage}%
}

\newcommand{\zvfrontmatter}{%
  \sloppy
  \makezvtitle
  \tableofcontents
  \newpage
}
